\subsection{Data and Methodology}
  \label{subsec:data_method}

  We use the SPARC database, which provides high-quality rotation curves together
  with homogeneous photometry and baryonic mass models for disk galaxies spanning
  a wide range of morphologies, surface brightnesses, and dynamical regimes~\cite{Lelli2016SPARC}.

  The numerical pipeline used to generate the rotation-curve figures and diagnostics is available as open-source
  software~\cite{CosmochronySimulation}.

  For each galaxy, the baryonic mass distribution is constructed from the observed stellar and gas components.
  Stellar mass-to-light ratios are fixed following the standard SPARC prescriptions.
  No additional dark matter component is introduced in this effective description.

  The effective acceleration $g_{\mathrm{eff}}(r)$ is obtained by solving
  Eq.~\eqref{eq:effective_acceleration} using the Newtonian acceleration
  $g_{\mathrm{N}}(r)$ computed from the baryonic mass distribution.
  The circular velocity then follows from
  \begin{equation}
    v_{\mathrm{eff}}(r) = \sqrt{r\, g_{\mathrm{eff}}(r)} .
  \end{equation}

  The saturation scale $a_\star$ is treated as a universal parameter.
  Its value is fixed globally by minimizing the total residuals across the
  sample, rather than optimized on a galaxy-by-galaxy basis.
